\subsection{Education Surrounding Algorithms}
Research surrounding the teaching of computer science concepts in secondary education has increasingly emphasized the importance of algorithmic
  thinking as a foundational component of computer science literacy \cite{futschek2006algorithmic}\cite{bell2014presenting}\cite{gal2004efficiency}.
Futschek argues that algorithmic thinking represents one of the central competencies of computer science education, encompassing abilities such as
  analyzing problems, identifying appropriate actions, constructing correct procedures and evaluating efficiency to support conceptual reasoning skills
  necessary to understand why algorithms function when they’ve likely encountered the relevant programming syntax \cite{futschek2006algorithmic}.

Ioannou and Angeli note that traditional computer science instruction has historically emphasized teacher-centered delivery and programming syntax at
  the expense of conceptual understanding.
They describe how secondary computer science education frequently ignored relationships between pedagogy, learner misconceptions, and subject matter
  understanding \cite{ioan2013}.
Their work instead calls for Technological Pedagogical Content Knowledge (TPCK), a framework that integrates pedagogy, content knowledge, learner
  understanding, and educational technology.
According to this framework, effective instruction requires teachers to identify abstract or difficult concepts and use representations, simulations,
  and interactive tools to make them accessible to secondary learners due to the inherent abstraction.

Work completed in New Zealand further suggests that algorithm instruction should extend beyond programming itself and include broader computational
  concepts.
Bell et al.
argue that high school computer science education often becomes overly focused on coding while neglecting theoretical ideas such as complexity, encryption, human-computer interaction, and algorithmic efficiency \cite{bell2014presenting}.
Their work emphasizes that students benefit from understanding how algorithms exist independently of programming languages and how different
  algorithms may solve the same problem with dramatically different performance characteristics.
This broader approach also helps students who may not initially find programming intrinsically motivating by connecting algorithms to practical and
  interdisciplinary applications.

Bell et al.
further recommends constructivist instructional methods for teaching algorithms effectively.
They describe instructional designs where students discover sorting and searching strategies through interactive games, and physical activities
  before formally studying the algorithms themselves \cite{bell2014presenting}.
Similarly, Futschek advocates teaching algorithmic thinking independently from programming through visual problem-solving activities such as maze
  navigation and pathfinding tasks \cite{futschek2006algorithmic}.
These approaches allow students to focus on reasoning and problem decomposition without simultaneously struggling with programming syntax.
Their studies suggest that this separation reduces mental load and enables students to build stronger conceptual models before implementation.

Literature additionally highlights the importance of visualization and interactive learning tools in algorithm instruction.
Visualization systems, simulations, and dynamic representations help students mentally model iterative and recursive processes that are otherwise
  difficult to observe directly.
Ioannou and Angeli emphasize that technologies capable of dynamic data transformation, simulation, and interactive representation are especially
  valuable for topics students traditionally struggle to conceptualize \cite{ioan2013}.
Bell et al.
similarly note that algorithm visualizations become most effective when students actively interact with them rather than passively observe them \cite{bell2014presenting}.

Research suggests that competitions and collaborative programming activities can support algorithmic learning and student motivation
  \cite{combefis2014programming}\cite{cuba2014learning}.
Research surrounding online and in-person programming competitions demonstrates that structured problem-solving environments encourage students to
  practice algorithmic reasoning repeatedly while receiving feedback on correctness and efficiency.
These environments also expose students to authentic computational challenges while helping develop persistence, debugging skills, and computational
  fluency.
However, researchers caution that contests alone are insufficient without instructional support, reflection, and conceptual discussion
  \cite{combefis2014programming}.
Effective algorithm education therefore appears to require a balance between theoretical understanding, guided practice, collaborative learning, and
  meaningful problem-solving experiences.

The clearest conflict from the above concerns the relationship between syntax and context.
Futschek’s studies differ in the degree to which they pedagogically prioritize conceptual reasoning before programming implementation compared to the
  works completed by Bell et al.
and Ioannu and Angeli \cite{futschek2006algorithmic}\cite{bell2014presenting}\cite{ioan2013}.
Though relatedly, research also disagrees on the role and efficacy of competition in learning algorithms.
Contest-based literature portrays programming competitions as highly motivating environments that improve algorithmic skills, persistence, and
  engagement \cite{combefis2014programming}\cite{cuba2014learning}.
However, Ioannou and Angeli, and Bell et al.
place greater emphasis on learner-centered pedagogy, accessibility, and conceptual understanding rather than competitive performance \cite{bell2014presenting}\cite{ioan2013}.
The implicit tension is that competition-based models may primarily benefit already high-performing or intrinsically motivated students, whereas
  constructivist and visualization-focused approaches are designed to support broader classroom populations, including beginners and less confident
  learners.

There is also disagreement regarding what aspects of algorithms should be emphasized at the secondary level.
Gal-Ezer and Zur focus heavily on efficiency and computational complexity, arguing that understanding efficiency is fundamental to algorithm
  education \cite{gal2004efficiency}.
Bell et al., however, cautions against an overly narrow focus on programming and efficiency alone, instead advocating for a broader computational
  literacy that includes areas like HCI, graphics, encryption, AI and others \cite{bell2014presenting}.
Additionally, while Futschek’s work emphasizes abstraction, modeling, and reasoning about algorithms conceptually through mazes and pathfinding tasks
  Bell et al.
disagree and argues that connecting student learning is better supported through a more hands on approach \cite{futschek2006algorithmic}\cite{bell2014presenting}.
Though they support conceptual understanding Futschek advocates they also stress connecting algorithms to authentic real-world computing contexts and
  interdisciplinary applications.
Contest literature, on the other hand, tends to prioritize optimization, performance, and solving constrained technical problems.
